\documentclass[10pt,a4paper]{amsart}
\usepackage{amsmath,amssymb,amsthm,graphicx}
\usepackage{verbatim}
\newenvironment{reference}{\comment}{\endcomment}
%\newenvironment{reference}{}{}
\newenvironment{slogan}{\comment}{\endcomment}
\newenvironment{history}{\comment}{\endcomment}

% For commutative diagrams we use Xy-pic
\usepackage[all]{xy}

% We use 2cell for 2-commutative diagrams.
\xyoption{2cell}
% Theorem environments.
%
\theoremstyle{plain}
\newtheorem{theorem}[subsection]{Theorem}
\newtheorem{proposition}[subsection]{Proposition}
\newtheorem{lemma}[subsection]{Lemma}

\theoremstyle{definition}
\newtheorem{definition}[subsection]{Definition}
\newtheorem{example}[subsection]{Example}
\newtheorem{exercise}[subsection]{Exercise}
\newtheorem{situation}[subsection]{Situation}

\theoremstyle{remark}
\newtheorem{remark}[subsection]{Remark}
\newtheorem{remarks}[subsection]{Remarks}

\numberwithin{equation}{subsection}

% Macros
%
\def\lim{\mathop{\mathrm{lim}}\nolimits}
\def\colim{\mathop{\mathrm{colim}}\nolimits}
\def\Spec{\mathop{\mathrm{Spec}}}
\def\Hom{\mathop{\mathrm{Hom}}\nolimits}
\def\Ext{\mathop{\mathrm{Ext}}\nolimits}
\def\SheafHom{\mathop{\mathcal{H}\!\mathit{om}}\nolimits}
\def\SheafExt{\mathop{\mathcal{E}\!\mathit{xt}}\nolimits}
\def\Sch{\mathit{Sch}}
\def\Mor{\mathop{\mathrm{Mor}}\nolimits}
\def\Ob{\mathop{\mathrm{Ob}}\nolimits}
\def\Sh{\mathop{\mathit{Sh}}\nolimits}
\def\NL{\mathop{N\!L}\nolimits}
\def\CH{\mathop{\mathrm{CH}}\nolimits}
\def\proetale{{pro\text{-}\acute{e}tale}}
\def\etale{{\acute{e}tale}}
\def\QCoh{\mathit{QCoh}}
\def\Ker{\mathop{\mathrm{Ker}}}
\def\Im{\mathop{\mathrm{Im}}}
\def\Coker{\mathop{\mathrm{Coker}}}
\def\Coim{\mathop{\mathrm{Coim}}}

% Boxtimes
%
\DeclareMathSymbol{\boxtimes}{\mathbin}{AMSa}{"02}

%
% Macros for moduli stacks/spaces
%
\def\QCohstack{\mathcal{QC}\!\mathit{oh}}
\def\Cohstack{\mathcal{C}\!\mathit{oh}}
\def\Spacesstack{\mathcal{S}\!\mathit{paces}}
\def\Quotfunctor{\mathrm{Quot}}
\def\Hilbfunctor{\mathrm{Hilb}}
\def\Curvesstack{\mathcal{C}\!\mathit{urves}}
\def\Polarizedstack{\mathcal{P}\!\mathit{olarized}}
\def\Complexesstack{\mathcal{C}\!\mathit{omplexes}}
% \Pic is the operator that assigns to X its picard group, usage \Pic(X)
% \Picardstack_{X/B} denotes the Picard stack of X over B
% \Picardfunctor_{X/B} denotes the Picard functor of X over B
\def\Pic{\mathop{\mathrm{Pic}}\nolimits}
\def\Picardstack{\mathcal{P}\!\mathit{ic}}
\def\Picardfunctor{\mathrm{Pic}}
\def\Deformationcategory{\mathcal{D}\!\mathit{ef}}
\title{An etale cover that cannot be refined by a finite composition of open immersions and finite etale morphisms}
\author{The Stacks Project Authors}
\date{stacks-project:examples.tex:local-snapshot-20261012; GNU FDL 1.2 or later}
\begin{document}
\maketitle
\noindent\textit{Editorial note.} All delivered mathematical text is the author's own, copied from the pinned Stacks Project chapter file examples.tex; no statement or proof step was rewritten, reordered or removed.\par\medskip
\noindent\textit{Editorial note.} The counted proof is the author's complete argument: Context span 4444-4452 included per the situation-above precedent; without it the statement's 'the etale cover U -> X' dangles.\par\medskip
\noindent\textit{Editorial note.} The label lemma-etale-not-dominated-zariski-finite-etale names the counted statement, and the proof environment immediately following it in the source is delivered verbatim as its proof.\par\medskip
\noindent\textit{Editorial note.} The counted statement is carried by the author's own lemma environment in the chapter source.\par\medskip
\noindent\textit{Allowed standard prerequisites (editorial).}\begin{itemize}
\item The section construction carried as context span: X a smooth connected affine curve over C, g : Y -> X finite of degree 3 with fibre over x equal to an unramified point u (degree 1) and a ramified point y (degree 2), U = Y \textbackslash{} \{y\} and f = g|\_U : U -> X etale surjective with f\textasciicircum{}\{-1\}(x) = \{u\}
\item Discrete valuation rings, ramification indices, connected components of smooth curves; fibres of finite morphisms
\end{itemize}\medskip
Let $X$ be a smooth, connected affine curve over $\mathbf{C}$.
Choose a finite morphism $g: Y \to X$ of degree $3$ from
a smooth connected curve $Y$ such that there is a point $x \in X$
whose fiber $g^{-1}(x) = \{u, y\}$ consists of an unramified point
$u$ (degree 1) and a ramified point $y$ (degree 2). By replacing
$X$ with a Zariski open neighborhood of $x$, we may assume $g$
is unramified everywhere except at $y$. Let $U = Y \setminus \{y\}$.
Then $f=g|_U: U \to X$ is an \'etale surjective morphism, and
$f^{-1}(x) = \{u\}$. 

\begin{lemma}
\label{lemma-etale-not-dominated-zariski-finite-etale}
The \'etale cover $U \to X$ cannot be refined by any finite composition
of open immersions and finite \'etale morphisms.
\end{lemma}

\begin{proof}
Towards contradiction, assume there exists a tower of morphisms 
$$
W_m \to W_{m-1} \to \dots \to W_1 \to W_0 = X
$$
where each map $W_i \to W_{i - 1}$ is either an open immersion or a
finite \'etale morphism, a lift $W_m \to U$ of $W_m \to X$ along
$f : U \to X$, and a point $w \in W_m$ mapping to $u \in U$
(and thus to $x \in X$). By replacing each $W_k$ with the connected
component containing the image $w_k \in W_k$ of $w$, we may assume
that each $W_k$ is smooth and connected.

\medskip\noindent
For each $0 \le i \le m$, consider the fiber product $Z_i = Y \times_X W_i$.
As $W_i \to X$ is \'etale and $Y$ is smooth, the curve $Z_i$ is a disjoint
union of smooth, connected curves and the maps $Z_i \to W_i$ are finite
of degree $3$ by base change. Clearly, the fibre of $Z_i \to W_i$
over $w_i$ consists of two points $z_i, v_i$ mapping to $y, u$ in $Y$.
The morphism $Z_i \to W_i$ is unramified at $v_i$ and ramified of degree $2$
at $z_i$. Let $n_i = \#\pi_0(Z_i)$. Since the fibre of $Z_i \to W_i$
has $2$ points we see that $n_i \in \{1, 2\}$. Moreover, $n_i$ can only
increase with $i$ as the transition maps $Z_i \to Z_{i - 1}$ are dominant. 
We clearly have $n_0 = 1$. Moreover, we must have $n_m \geq 2$: there is an
$X$-morphism $W_m \to U \subset Y$, so the base change $Z_m \to W_m$
has a section. There is then a minimal index $i \ge 1$ where $n_{i - 1} = 1$
and $n_i = 2$. As the $n_j$'s are determined by whether or not $Z_i$
is irreducible, the map $W_i \to W_{i - 1}$ cannot be an open immersion,
so it must be finite \'etale. As $n_i = 2$, the smooth curve $Z_i$
has two connected components; as the map to $W_i$ is finite, one of
these components must then map isomorphically to $W_i$,
so we can write $Z_i = W_i \sqcup Z'_i$, where $Z'_i \to W_i$ has degree $2$.
The induced map $W_i \subset Z_i \to Z_{i - 1}$ is a dominant map
between connected curves, both of which are finite over $W_{i - 1}$.
Any such map must be surjective. Thus we may pick a point $t \in W_i$
mapping to $z_{i - 1}$. We obtain extensions 
$$
\mathcal{O}_{W_{i - 1}, w_{i - 1}} \to
\mathcal{O}_{Z_i, z_i} \to
\mathcal{O}_{W_i, z}
$$
of discrete valuation rings. Now the first one is ramified while
the composite is unramified as $W_i \to W_{i - 1}$ was \'etale; this is
impossible, so we obtain a contradiction.
\end{proof}

\end{document}
